\documentclass[12pt]{article}
\usepackage{personal}
\usepackage{realoracles}


\newtheorem{lemma}{Lemma}[section]
\newtheorem{corollary}{Corollary}[section]
\newtheorem{proposition}{Proposition}[section]
\newtheorem{theorem}{Theorem}


\newcommand{\qcut}[2][x]{\ensuremath{\mathbb{Q}_{#2 #1}}}
\newcommand{\qlt}[1][x]{\qcut[#1]{<}}
\newcommand{\qeq}[1][x]{\qcut[#1]{=}}
\newcommand{\qgt}[1][x]{\qcut[#1]{>}}
\newcommand{\qgeq}[1][x]{\qcut[#1]{\geq}}
\newcommand{\qleq}[1][x]{\qcut[#1]{\leq}}
\newcommand{\cut}[1][x]{{\vert}_{#1} }

\newcommand{\yrel}{\xrel[y]}

\title{Real Numbers, Rational Betweenness Relations, and Dedekind Cuts}
\jtauthor
\date{\today}

%\sloppy%\openup-.1\jot
\begin{document}\maketitle
\begin{abstract}
How well known a real number is can be specified by a rational interval that contains the real number. Flipping this leads to defining a real number by a rational betweenness relation. Specifically, two rational numbers are $x$-related if $x$ ought to be, inclusively, between them. These betweenness relations represent a new way to define the real numbers. There are five properties that characterize such relations. The key property is the Separation property which states that given two rationals in relation, then any number in between is either $x$-related to exactly one of the two endpoints or is $x$-related to itself. To prove that the set of such relations can be taken to be the real numbers, it will be demonstrated that the betweenness relations with the natural order and arithmetic of intervals is canonically isomorphic as a complete ordered field to Dedekind cuts. 
\end{abstract}

\section{Introduction}

The purpose of this paper is to give a new definition of real numbers and establish that this definition works. 

The first task is to define the term rational betweenness relation. After establishing the definition and some useful properties, the second task is to give the definition of Dedekind cut to use and its relevant properties. The first main result is to establish the explicit canonical bijection from the space of Dedekind cuts to the space of rational betweenness relations. 

An important part of defining real numbers is defining the order relation and arithmetic operators. Interval ordering and arithmetic naturally carry over to the relations. The bijection to the Dedekind cuts will be shown to respect the order relations. For the  arithmetic operators, the bijection is used to define them, but in a way which aligns with what would seem to be the natural way to use the arithmetic of intervals in doing so.  The bijection is then used to establish the real number properties, including the completeness property.   This establishes that the space of rational betweenness relations is a valid model of the real numbers. 

Throughout this paper and generally without further explicit comment, the letters $a, b, c, d$ will always be rational numbers and often take the role of endpoints of intervals. The letters $m, n, p, q, r, s$ will also always be rational numbers and often take the role of numbers in an interval. The symbols $\varepsilon$ and $\delta$ will be non-negative rational numbers. The letters $x, y, z, \alpha, \beta$ will notionally represent real numbers.

 A \textbf{rational interval} $a \rint b$ is the set of all rationals $q$ such that $q$ is between $a$ and $b$, inclusively. A \textbf{rational singleton} $a:a$ is a set consisting of just $a$ and it is accepted as a rational interval. The empty set is not a rational interval. 

Similar notation will be used to indicate betweenness such as $a \rint b\rint c$ implying $b$ is, inclusively, between $a$ and $c$. It is easy to extend this notation such as $a \rint b \rint c \rint d$ implying either $a \leq b \leq c \leq d$ or $d \leq c \leq b \leq a$. Replacing $\rint$ with $\betneq$ implies strict inequality. A bracketed expression in a chain such as $\{b, c\}$ is used to indicate common betweenness relations around $b$ and $c$ without deciding the relative ordering of $b$ and $c$. Thus, $a \betneq \{b, c \} \rint d$ would be interpreted as $a$ is not equal to $b$ or $c$ and that both $b$ and $c$ are between $a$ and $d$. 

If $b$ is not between $a$ and $c$, this could be indicated either by \sout{$a \rint b\rint c$} or by $b\rint \{a,c\}$.

\section{Definition of a Rational Betweenness Relation}

A rational betweenness relation is a relation on pairs of rational numbers that is supposed to express that there is a real number $x$ between the two rational numbers. The notation $a \xrel b$ denotes that $x$ ought to satisfy $a \leq x \leq b$ or $b \leq x \leq a$;  both the ordered pairs $(a,b)$ and $(b,a)$ are then in the betweenness relation. The interval $a \rint b$ is then an \textbf{$x$-interval} and $a$ is \textbf{$x$-related} to $b$. If $(a,b)$ is not in the relation, then that can be denoted by \sout{$a \xrel b$}. 

Since this is defining the real number $x$, the above is only a notional idea behind the definition. The definition is that a \textbf{rational betweenness relation} is a symmetric relation on rational numbers 
which satisfies the following properties:
\begin{enumerate}
    \item Existence. There exists $a$ and $b$ such that $a\xrel b$.
    \item Interval Separation. If $a \xrel b$ and $a \rint  c \rint  b$, then exactly one of the following holds: 1) $a \xrel c$ and \sout{$c \xrel b$}, 2) $c \xrel b$ and \sout{$a \xrel c$}, or 3) $c \xrel c$. 
    \item Consistency. If $c \rint  a \rint b \rint  d$ and $a \xrel b$, then $c \xrel d$. 
    \item Singular. If $c \xrel c$ and $d \xrel d$, then $c=d$. 
    \item Closed. If $c$ is a rational number such that $c$ is included in every $x$-interval $a \rint b$, then  $c \xrel c$. 
\end{enumerate}

If $q \xrel q$, then $x$ represents the rational number $q$ and $q$ is the \textbf{root of the relation}.
If such a $q$ exists, then $x$ is said to be rational. Also note that the Singular property implies that $q$ would be unique. If no such $q$ exists, then $x$ is said to be irrational. 

The Interval Separation property is at the heart of this procedure being computationally useful. It is modeled on how the Intermediate Value Theorem is often operationalized. For example, if $f$ is strictly monotonic, defined on all rational numbers, and has both positive and negative values, then the rational betweenness relation $x$ representing the location where the sign changes is defined as $a \xrel b$ exactly when $f(a)*f(b) \leq 0$. To describe the positive $n$-th root of $q$, the relation is $a \xrel b$ for $a \leq b$, exactly when $\max(0, a)^n \rint q \rint b^n$. For a Cauchy sequence, the relation can be described prosaically as the intervals that contain the tail of the sequence. 

\section{Useful Statements for Rational Betweenness Relations}

Here are some useful general statements about rational betweenness relations. 

\begin{lemma}\label{br--rooted}
    If the $x$-between relation has a rational root $q$, then $a \xrel b$ if and only if $a \rint q \rint b$.
\end{lemma}

Note that $a \xrel b$ is stating that the interval $a \rint b$ is in the relation while $a \rint q \rint b$ is stating that the rational number $q$ is between the rational numbers $a$ and $b$.

\begin{proof}
    Let $a \rint b$ be given such that $a \rint q \rint b$ which is the same as $a\rint q\rint q\rint b$. Since $q \xrel q$, Consistency implies that $a \xrel b$. For the other direction, assume that \sout{$a \rint q \rint b$} which implies $q$ is distinct from both $a$ and $b$. By Trichotomy and without loss of generality, assume that $a$ is between $q$ and $b$. By Consistency using $q \xrel q$, $q \xrel b$ and $q\xrel a$.  Singular implies that \sout{$a \xrel a$}. Since $q \xrel a$ and $q \rint a \rint b$, Separation, using $a$ as the Separation point, states that \sout{$ a \xrel b$}. Since either $a \rint q \rint b$ or \sout{$a \rint q \rint b$} must hold in general, this has exhausted the cases and the result holds.  
\end{proof}

Every rational is $x$-related to some other rational for every betweenness relation $x$. 

\begin{lemma}\label{br--existence}
    Given $\xrel$ and rational $a$, there exists $b$ such that $a \xrel b$. 
\end{lemma}

\begin{proof}
    Let $u \xrel v$ be an $x$-interval; that such an interval exists follows from the Existence property. If $a$ is not between $u$ and $v$, then without loss of generality, it can be assumed that $a\rint u\rint v\rint v$. Then Consistency tells us that $a \xrel v$. In this case, $v$ is the required $b$. 

    The other case is $u\rint a\rint v$. Apply Separation, using $a$ as the Separation point, to obtain that either 1) $u \xrel a$, 2) $a \xrel v$, or $a \xrel a$. In all three cases, there exists the desired $b$. 
\end{proof}

Two $x$-related intervals that intersect just at an endpoint leads to that endpoint being the root of $x$. 

\begin{proposition}\label{br--endpointed}
    If $a \xrel p$ and $p \xrel b$ with $a\rint p\rint b$, then $p \xrel p$.
\end{proposition}

\begin{proof}
    By Consistency, $a \xrel b$ as $a\rint a\rint p\rint b$.  By Separation, using $p$ as the Separation point, exactly one of the following holds true: $p \xrel p$,  \sout{$a \xrel p$}, or \sout{$p \xrel b$}. Since both $a\rint p$ and $p\rint b$ are in the relation by assumption, it must be the case that $p \xrel p$, as was to be shown. 
\end{proof}

All $x$-intervals overlap and their intersection is itself an $x$-interval. This makes it possible to condense the information from multiple $x$-intervals. 

\begin{proposition}\label{br--intersect}
    If $a \xrel b$ and $c \xrel d$, then $a \rint b$ and $c\rint d$ intersect in an $x$-interval. If it is also the case that  $\{a, c\} \rint \{b, d\}$, then  $c \xrel b$ and $a \xrel d$.
\end{proposition}

\begin{proof}
    The betweenness of the four rational numbers $a$, $b$, $c$, $d$ is well handled by cases. To minimize the cases, without loss of generality, assume that $a$ is one of the outer endpoints and $c$ is closer in the betweenness string to $a$ than $d$ is. We have the following cases: 

    \begin{enumerate}
        \item $a \rint b \betneq c\rint d$. Let $m = \frac{b+c}{2}$. This is a rational number strictly between $b$ and $c$ with the strictness following from them not being equal to each other. Thus, $a \rint b\rint m\rint c\rint d$. By Consistency, $a\xrel m$ and $m \xrel d$. From Proposition \ref{br--endpointed}, this implies $m \xrel m$. But then lemma \ref{br--rooted} states that since $m$ is not in $a \rint b$ nor $c\rint d$, they cannot be $x$-intervals. Having arrived at a contradiction, this case is incompatible with the assumptions of the statement. 
        \item $a\rint (b=c)\rint d$. By Proposition \ref{br--endpointed}, $b\rint b$ is an $x$-singleton. This is the intersection of the two intervals and it is an $x$-interval. Note that $c \xrel b$ and, as $a \rint d$ contains it, by Consistency $a \xrel d$.
        \item $a\rint c\rint d\rint b$. The interval $c\rint d$ is the intersection and it was given as an $x$-interval. By Consistency, both $a \xrel d$ and $b \xrel c$.
        \item $a\rint c\betneq b\rint d$. The interval $c\rint b$ is the intersection. By applying Separation, with $c$ as the Separation point in the interval $a \rint b$, one of the following holds true: 
        
\begin{enumerate}
            \item $c \xrel b$. Since $c\rint b$ is the intersection, this case is exactly what is desired.
            \item $c \xrel c$. Consistency then implies $c \xrel b$ which is the intersection. 
            \item $a \xrel c$. Since $c \xrel d$ and $a\rint c\rint d$, Proposition \ref{br--endpointed} implies $c \xrel c$. The previous case then applies. 
        \end{enumerate}
        As $b \xrel c$ and $a \rint d$ contains $b \rint c$, Consistency yields $a \xrel d$.
    \end{enumerate}
\end{proof}


\begin{corollary}\label{br--cor-disj}
    If $a \xrel b$ and $c \rint d$ is disjoint from $a \rint b$, then \sout{$c \xrel d$}. 
\end{corollary}

The following is a useful restatement of the Closed property. 

\begin{lemma}\label{br--not-rooted-separated}
    If \sout{$a \xrel a$}, then there exists an $x$-interval $u\rint v$ that does not contain $a$. 
\end{lemma}

\begin{proof}
    The set of all $x$-intervals is not empty by the Existence property. If all $x$-intervals contained $a$, then the Closed property would imply that $a \xrel a$. Since  \sout{$a \xrel a$}, there must be some $x$-interval that does not contain $a$. 
\end{proof}

The next two statements suggest a narrowing away from endpoints that are not roots of $x$. 

\begin{lemma}\label{br--something-inside}
    If $a \xrel b$, then either $a \xrel a$ or there exists $c$ such that $a \betneq  c \rint  b$ with $c \xrel b$. 
\end{lemma}

\begin{proof}
    Assume \sout{$a \xrel a$}. By Lemma \ref{br--not-rooted-separated}, there exists an $x$-interval $u\rint v$ which does not contain $a$. By Proposition \ref{br--intersect}, there exists an $x$-interval  $c\rint d$ which is the intersection of $u\rint v$ and $a \rint b$. The intersection interval excludes $a$ as $a$ is not in $u\rint v$. Without loss of generality, it can be assumed that $a\betneq c\rint \{b,d\}$ which then yields $c \xrel b$.
\end{proof}

\begin{corollary}\label{br--doublesomething-inside}
    If $a \xrel b$ and neither $a$ nor $b$ are roots of $x$, then there exists $c \rint d$ such that $a \betneq  c \rint  d \betneq  b$ with $c \xrel d$. 
\end{corollary}

\begin{proof}
Apply Lemma \ref{br--something-inside} to obtain $c \xrel b$. Then apply the lemma again to endpoint $b$ of $c \xrel b$ to obtain $c \xrel d$ satisfying $a \betneq c \rint d \betneq b$.
\end{proof}

Disjointness is a useful condition to use in separating $x$-intervals from non-$x$-intervals. 

\begin{lemma}\label{br--something-outside}
    If \sout{$a \xrel b$}, then there exists an $x$-interval disjoint from $a \rint b$.
\end{lemma}

\begin{proof}
    Since \sout{$a \xrel b$}, all subintervals, including the singletons $a\rint a$ and $b\rint b$, are not $x$-intervals by Consistency. From Lemma \ref{br--not-rooted-separated}, there exist $x$-intervals $A\rint A'$ and $B\rint B'$ such that $a$ and $b$ are not in those intervals, respectively. Since $x$-intervals intersect in an $x$-interval, there exists an $x$-interval $c\rint d$ which excludes $a$ and $b$. Thus, $c\rint d$ is either disjoint from $a \rint b$ or contained in $a \rint b$. It cannot be contained in $a \rint b$ as Consistency would then imply $a \rint b$ is an $x$-interval. Thus, the $x$-interval $c\rint d$ is disjoint from $a \rint b$.
\end{proof}


Two rational betweenness relations $\xrel$ and $\xrel[y]$ are different if there exists at least one interval $a \rint b$ such that either $a \xrel b$ with \sout{$a \xrel[y] b$} or $a \xrel[y] b$ with \sout{$a \xrel b$}. 

\begin{proposition}\label{br--different}
    If $\xrel$ and $\xrel[y]$ are different, then there exists two disjoint intervals $a \rint b$ and $c\rint d$ such that $a \xrel b$, \sout{$a \xrel[y] b$}, $c \xrel[y] d$, and \sout{$c \xrel d$}.
\end{proposition}


\begin{proof}
    By the definition of difference there exists an interval which is related for the one relation, but not the other; without loss of generality, assume that the interval $a \rint b$ is an $x$-interval but not a $y$-interval. By Lemma \ref{br--something-outside}, there exists a $y$-interval $c\rint d$ disjoint from $a \rint b$. Since it is disjoint from $a \rint b$, Corollary \ref{br--cor-disj} implies $c\rint d$ is not an $x$-interval. This is what was to be shown. 
\end{proof}


The final proposition demonstrates that there are arbitrarily small intervals in the relation. 

\begin{proposition}
    Given rational $\varepsilon >0$ and rational betweenness relation $\xrel$, there exists $a \rint b$ such that $|a-b| < \varepsilon$ and $a \xrel b$. 
\end{proposition}

This is the bisection algorithm as used, for example, in some proofs of the Intermediate Value Theorem.

\begin{proof}
    Let $a_0\rint b_0$ be an $x$-interval; this exists by the Existence property.
    
    Assume $a_i\rint b_i$ is defined. Let $m_i = \frac{a_i + b_i}{2}$. By Separation, using $m_i$ as the Separation point in $a_i \rint b_i$, either 1) $a_i \xrel m_i$, 2) $b_i \xrel m_i$, or 3) $m_i \xrel m_i$. If it is 3), then $m_i\rint m_i$ satisfies the requirements as its length is 0. Otherwise, for 1), let $a_{i+1}= a_i$, $b_{i+1} = m_i$ and for 2) let $a_{i+1} = m_i$, $b_{i+1} = b_i$. Then $|a_{i+1} - b_{i+1}| = 2^{-1} |a_i-b_i|= 2^{-i} |a_0-b_0|$. 

    Since there exists an $N$ such that $2^{-N} |a_0-b_0| < \varepsilon$, this process will finish in at most $N$ steps. 
\end{proof}




\section{Dedekind cuts}

A Dedekind cut models real numbers as a dividing point between rationals. There are multiple definitions. The one used here is a symmetric definition. 

A Dedekind cut is defined as a triplet of sets of rationals, $(A, B, C)$, such that: 
\begin{enumerate}
    \item Nonempty. Both $A$ and $C$ are non-empty.
    \item Ordered. If $a \in A$ and $c \in C$ then $a < c$. If $b \in B$, then also $a < b < c$. 
    \item Comprehensive. $A \cup B \cup C = \mathbb{Q}$.
    \item Singular. There is at most one element in $B$.
    \item Open. There is no greatest element in $A$ and no least element in $C$.
\end{enumerate}
If the cut ought to represent $x$, then the notation that will be used will be $\qlt$ for the set $A$, $\qeq$ for the, possibly empty, set $B$, and the set $\qgt$ for $C$. The notation $\qleq$ will be used for the set $\qlt \cup \qeq$ while $\qgeq$ denotes $\qgt \cup \qeq$. To refer to the cut itself, the notation $\cut$ will be used. 



If $\qeq = \{q\}$, then $\qlt$ is the set of rational numbers less than $q$ and $\qgt$ is the set of rational numbers greater than $q$. The notation $\cut[q]$ will represent this cut, noting that $q$ is a rational number. Cuts of this form represent the rational numbers and $q$ is the root of the cut. 

If $\qeq$ is empty, then the cut represents an irrational number. In that case, $\qgt$ and $\qlt$ are complements in the set of rational numbers. 

A common definition of Dedekind cuts is to define them solely in terms of the sets $\qlt$. This is what is done in \cite{rudin}, page 17. One can go from either definition to the other as established below. 


\section{Useful Statements for Dedekind Cuts}

This is a collection of some useful statements regarding Dedekind cuts. 

\begin{lemma}
    For a cut $x$, the sets $\qlt$, $\qeq$, $\qgt$ are pairwise disjoint. 
\end{lemma}

\begin{proof}
    This follows from the Ordered property. If an element was in two of the sets, then it would be less than itself which cannot be. 
\end{proof}

The lower set $\qlt$ is unbounded below while the upper set $\qgt$ is unbounded above. 

\begin{lemma}\label{br--cut-unbounded}
    Let $x$ be a cut, then the following hold true:
    \begin{enumerate}
        \item If $a \in \qleq$ and $p < a$, then $p \in \qlt$. 
        \item If $b \in \qgeq$ and $b < q$, then $q \in \qgt$. 
        \item If $a \in \qleq$ and $p \leq a$, then $p \in \qleq$. 
        \item If $b \in \qgeq$ and $b \leq q$, then $q \in \qgeq$. 
    \end{enumerate}
\end{lemma}

\begin{proof}
    For item 1, let $a$ and $p$ be as in the statement. By the Comprehensive property, $p$ is an element of either $\qlt$ or $\qgeq$. If it is in $\qgeq$, then $p \geq a$ by the Ordered property, but it is assumed that $p < a$. Thus, $p \in \qlt$.

    Item 2 follows similarly, using that if $q \in \qleq$, then $q \leq b$ which it is not. 

    For item 3, if $p = a$, then $p \in \qleq$. If $p < a$, then item 1 applies. 

    Similarly, for item 4, if $q = b$, then $q \in \qgeq$. If $b < q$, then item 2 applies. 

\end{proof}

\begin{lemma}\label{br--cut-equal}
    Two cuts $x$ and $y$ are equal if $\qlt \subseteq \qlt[y]$ and $\qgt \subseteq \qgt[y]$.
\end{lemma}


\begin{proof}
    By Comprehension and Disjointness, any given $q$ is in exactly one of $\qlt$, $\qeq$, or $\qgt$ and is in exactly one of $\qlt[y]$, $\qeq[y]$, or $\qgt[y]$. Since $q \in \qlt$ implies $q \in \qlt[y]$ and $q \in \qgt$ implies $q \in \qgt[y]$, the only possible case to question is $q \in \qeq$. By Singular, there is at most one $q \in \qeq$. Assume $q \in \qeq$ but that $q \notin \qeq[y]$. Then $q$ is either in $\qlt[y]$ or $\qgt[y]$. If $q \in \qlt[y]$, then, by the Open property, there exists $r > q$ such that $r \in \qlt[y]$. But $r > q$ implies $r \in \qgt $ by the previous result as $q \in \qgeq$. That would, however, imply $r \in \qgt[y]$ which contradicts the Disjointness of $\qlt[y]$ and $\qgt[y]$. Thus, $q \notin \qlt[y]$. Similarly, $q \in \qgt[y]$ would lead to an $r  < q$ such that $r \in \qgt[y]$ and $r \in \qlt$ implying $r \in \qlt[y]$ which would be a contradiction. Thus, $q \in \qeq[y]$ if $q \in \qeq$. As the other rationals were already handled by the hypothesis, it is the case that every rational belongs to the same corresponding component of $\cut$ and $\cut[y]$ which implies the cuts are equal.
    
\end{proof}

A set of rationals $A$ is a \textbf{lower cut} if it satisfies the following three properties: 
\begin{enumerate}
    \item Not Trivial. $A \neq \emptyset$ and $A \neq \mathbb{Q}$.
    \item Unbounded Below. If $p \in A$ and $q < p$, then $q \in A$. 
    \item Open Above. If $p \in A$, then there exists $q \in A$ such that $p < q$.
\end{enumerate}
The notation $A^c$ will denote the set complement of $A$. That is, $A$ and $A^c$ are disjoint with $A \cup A^c = \mathbb{Q}$. $A$ not being all of the rationals implies that $A^c$ is not empty. 

\begin{lemma}
    If $A$ is a lower cut, then $A^c$ is exactly the set of upper bounds of $A$. 
\end{lemma}

\begin{proof}
    Let $q \in A$ and $r \in A^c$. Since the sets are disjoint, $r \neq q$. If $r < q$, then $r \in A$ by the Unbounded Below property. This contradiction leads to the conclusion that $r \nless q$. By the Trichotomy property for rationals, $r > q$. As $q$ and $r$ were arbitrary elements in their sets, all of the elements in $A^c$ are strictly greater than the elements in $A$. 
    
    To establish that all upper bounds are in $A^c$, let $s$ be an upper bound which means that for any $t \in A$, $s \geq t$. If $s \in A$, then, by the Open Above property, there would be an $m \in A$ such that $m > s$. But this contradicts $s$ being an upper bound of $A$. Thus, we can conclude that $s \in A^c$. 
\end{proof}

\begin{corollary}
    The set $A^c$ is unbounded above and is not finite.     
\end{corollary}

\begin{proof}
    Since $A^c$ is nonempty, let $r \in A^c$. Let $s$ be any positive rational implying $r + s > r$. Using the lemma, for any $q \in A$, $r > q$ implying $r+s > q$.  As $q$ was an arbitrary element of $A$, this implies $r+s$ is an upper bound of $A$ and thus, by the lemma, is in $A^c$.
\end{proof}

\begin{proposition}\label{br--lowcutdef}
    If a set $A$ is a lower cut, then there is a unique cut $\cut$ such that $A = \qlt$.
\end{proposition}

An overline will be used to indicate the completion of the lower cut to a full cut. In the above, $\overline{A} = \cut$.

\begin{proof}
    Let $A^c$ be the complement of $A$. If there exists $q \in A^c$ such that $q \leq r$ for all $r \in A^c$, then define $B = \{q\}$ and $C = A^c \setminus \{q\}$. If there is no such least number $q$, then define $B = \{\}$ and $C = A^c$. 

    The claim is that $(A, B, C)$ forms a cut. The properties hold as follows:

    \begin{enumerate}
        \item Nonempty. The lower cut $A$ is nonempty by definition of a lower cut. The set $A^c$ is not finite by the above corollary. Thus, $A^c \setminus \{q\}$ would also not be finite and nonempty. Therefore, $C$ is nonempty. 
        \item Ordered. Let $a \in A$ and $c \in C$. Then the lemma says that $ a< c$. If $b \in B$, then the lemma also says that $a < b$. As $b$ is the least element of $A^c$ and $b \notin C$, $ b < c$. 
        \item Comprehensive. $A \cup B \cup C = A \cup A^c = \mathbb{Q}$ by definition. 
        \item Singular. The set $B$ consists of the least elements of $A^c$ and it is the case that this is either an empty set or singular set. Indeed, given any set $D$ of rationals, if $q, s \in D$, $q \leq r$ and $s \leq r$ for all $r \in D$, then $q \leq s$ and $s \leq q$ which implies $q = s$. That is, there is at most one least element for any given set. 
        \item Open. The Open Above property from $A$ being a lower cut is exactly that $A$ has no greatest element in $A$. If $A^c$ does not have a least element, then $A^c = C$ and $C$ does not have a least element. If $A^c$ does have a least element, say $q$, then the claim is that $C= A^c - {q}$ has no least element. By the lemma, $A^c$ is the set of upper bounds of $A$. In particular, since $q \in A^c$, it is an upper bound. Given any number $r \in C$, $r > q$ as $q$ is the least element of $A^c$. Since $m = \frac{q+r}{2}$ is between $r$ and $q$, it is the case that $r > m > q > t$ for all $t \in A$. Thus, $m$ is an upper bound of $A$ and therefore $m \in A^c$. Since $m$ is not equal to $q$, $m \in C$. As $m$ is less than $r$ and $r$ was an arbitrary element of $C$, there is no least element of $C$. 
    \end{enumerate}

    As for uniqueness, any cut such that $A = \qlt$ will have that $A^c = \qgeq$. This set will either have a least element or it will not and that determines whether $\qeq$ is singular or empty, respectively. Either way, $\qeq$ and $\qgt$ are uniquely determined by $\qlt$. 
\end{proof}


\section{From Relation to Cut and Back Again}

This section will establish a canonical bijection between the rational betweenness relations and Dedekind cuts. The term canonical refers to the fact that the natural bijection described below does not require any arbitrary choices. It maps the representative of rationals to each other, e.g., the zero cut is mapped to the zero relation. 

Given the betweenness relation $\xrel$, the associated Dedekind cut is defined by \begin{enumerate}
    \item $\qlt = \{ a | \exists\, b, c \, (a < c < b) \wedge (c \xrel b) \} $
    \item $\qgt = \{ b | \exists\, a, c \, (a < c < b) \wedge (a \xrel c) \} $
    \item $\qeq = \{ q | q \xrel q \} $
\end{enumerate}
Essentially, $\qlt$ is the set of non-root lower endpoints of $x$-intervals, $\mathbb{Q}_{>x}$ is the set of non-root upper endpoints of $x$-intervals, and $\qeq$ is the set of endpoints that form singletons in the $x$-between relation. The use of $c$ is present to ensure that if a singleton is present in the relation, then its endpoint does not appear in either of the inequality sets. The set $\qeq$ consists exactly of the roots of the relation; the Singular property asserts that there is at most one root. 

Is this a cut? Yes: 
\begin{enumerate}
    \item Nonempty. Existence yields an $x$-interval $a \rint b$. Without loss of generality, assume $a \leq b$. Then $a-1< a$, $b+1 > b$. By definition of the cut, $a-1 \in \qlt$ and $b+1 \in \qgt$.
    \item Ordered. Let $a \in \qlt$ and $b \in \qgt$.  Then there exists $c, m$ such that $a < c < m$ and $c \xrel m$. There is also a $d, n$ such that $d < n< b$  and $d \xrel n$. By Proposition \ref{br--intersect}, $c \xrel m$ and $d \xrel n$ intersect. Let $p$ be an element of that intersection. Then $a < c \leq p \leq n < b$ implying $a < b$.  If $\qeq$ is not empty, then let $q$ be the element of $\qeq$. By Lemma \ref{br--rooted}, $c \xrel m$ and $d \xrel n$ implies that both $a\rint c\rint q\rint m$ and $d\rint q\rint n\rint b$ hold true. Thus, $a < q$ and $q < b$.  
    \item Comprehensive. Given any rational $a$, Lemma \ref{br--existence} yields the existence of $b$ such that $a \xrel b$. Lemma \ref{br--something-inside} establishes the existence of a $c$ such that $a\betneq c\rint b$ with $c \xrel b$ unless $a \xrel a$ in which case $a \in \qeq$. In the other case, $a \in \qlt$ if $a < c$ while $a  \in \qgt$ if $a > c$. 
    \item Singular. By the Singular property of the relation, there is at most one $q$ such that $q \xrel q$. Thus $\qeq$ is either empty or consists of one element. 
    \item Open. Let $a\in \qlt$ and $b \in \qgt$. The task is to establish that there exists $m \in \qlt, n \in \qgt$ such that $a < m$ and $n < b$. By assumption, there exists an $x$-interval that excludes $a$ from above and an $x$-interval that excludes $b$ from below. By Proposition \ref{br--intersect}, those intervals intersect in an $x$-interval which can be taken to be $c\rint d$ with $c \leq d$. This implies $a\rint c\rint d\rint b$ with $c$ distinct from $a$ and $d$ distinct from $b$. Since $c$ or $d$ could be roots of the relation, it is necessary to look at other elements. Let $m = \frac{a+c}{2}$ and $n = \frac{d+b}{2}$. Then $m \in \qlt$ and $n \in \qgt$ as $c$ and $d$, respectively, still fill the same role as required in the definition. Thus $\qlt$ has no greatest element and $\qgt$ has no least element. 
\end{enumerate}

In the other direction, given the Dedekind $x$-cut, the associated $x$-betweenness relation is defined by, for $a \leq b$, $a \xrel b$ exactly when $a \in \qleq$ and $b \in \qgeq$. 

Is this a betweenness relation? Yes: 
\begin{enumerate}
    \item Existence. By the Nonempty property, we have the existence of $a \in \qlt$ and $b \in \qgt$ which implies $a \xrel b$. 
    \item Separation. Let $a \xrel b$ with $a \leq b$.  Then $a \in \qleq$ and $b \in \qgeq$. Let $c$ be given such that $a\rint c\rint b$. By the Comprehensive property, there are three possible cases: $c \in \qlt$, $c \in \qgt$, or $c \in \qeq$. If $c \in \qlt$, then  $c \xrel b$. It is also the case that \sout{$a \xrel c$} as both $a$ and $c$ are in $\qlt$ implying neither of them are in $\qgeq$. Similarly, if $c \in \qgt$, then \sout{$c \xrel b$} and $a \xrel c$. If $c \in \qeq$, then $c \in \qleq$ and $c \in \qgeq$ which implies $c \xrel c$. 
    \item Consistency. Let $c\rint a \rint b\rint d$ with $ a \xrel b$. The claim is that $c \xrel d$. Without loss of generality, assume $a \leq b$ which implies $a \in \qleq$ and $b \in \qgeq$.  The betweenness leads to $c \leq a$ and $ b \leq d$. By Lemma \ref{br--cut-unbounded},  $c \in \qleq$ and $ d \in \qgeq$. Thus, $c \xrel d$.
    \item Singular. The three pieces of the cut are disjoint. The only way to get an $x$-singleton $q\rint q$ is if $q \in \qeq$. That set has at most one element. 
    \item Closed. Assume that $c$ is such that, for any $a$, $b$, if $a \xrel b$, then $a\rint c\rint b$. Since the three cut sets disjointly span the rationals by the Comprehensive property, $c$ is in exactly one of those sets. If $c \in \qlt$, then, by the Open property, there exists $a \in \qlt$ such that $c < a$. Take any $b \in \qgt$ which can be done by the Nonempty property. Then $a \xrel b$ by definition while \sout{$a\rint c\rint b$} contradicting the assumption about $c$. Similarly, if $c \in \qgt$, then there would exist $b \in \qgt$ with $b < c$ and any $a \in \qlt$ would lead to \sout{$a\rint c\rint b$} with $a \xrel b$. With neither of these options available, it must be the case that $c \in \qeq$ implying $c \xrel c$.    
\end{enumerate}

The final task is to establish that these two are bijections and, in particular, they are the inverses of one another. 

To show the mapping from relations to cuts is one-to-one, let $\xrel$ and $\yrel$ represent two distinct rational betweenness relations. By Proposition \ref{br--different}, there exists $x$-interval $a \rint b$ and $y$-interval $c\rint d$ such that they are disjoint. Without loss of generality, assume that $a \leq b < c \leq d$. Then $c \in \qleq[y]$, $d \in \qgeq[y]$ while $c, d \in \qgt[x]$. Thus, different relations are mapped to different cuts. 

To show the mapping from the relations to the cuts is onto, the path will be to show that the cut-to-relation mapping is undone by the relation-to-cut mapping. Let a cut $\cut$ be given. Let $\yrel$ represent the relation that the cut is mapped to. Let $\cut[z]$ represent the cut that $\yrel$ is mapped to. The goal is to show that $\cut = \cut[z]$.  

Let $a \in \qlt[z]$ and $b \in \qgt[z]$. The main task is to show $a \in \qlt$ and $b \in \qgt$. By the definition of the relation-to-cut mapping, there exists $m_a < n_a, m_b < n_b$ such that $m_a \xrel[y] n_a$, $a < m_a$, and $m_b \xrel[y] n_b$ with $n_b < b$. The two $x$-intervals intersect in an $x$-interval $c\rint d$ by  Proposition \ref{br--intersect}. Without loss of generality, assume $c \leq d$. By transitivity, it is then the case that $a < c \leq d < b$ and $c \yrel d$. By the definition of the cut-to-relation mapping, $c \in \qleq$ and $d \in \qgeq$. Because $a < c$ and $ d < b$, $a \in \qlt$ and $b \in \qgt$ by Lemma \ref{br--cut-unbounded}. Thus, $\qlt[z] \subseteq \qlt$ and $\qgt[z] \subseteq \qgt$. By Lemma \ref{br--cut-equal}, $\cut[z] = \cut$

% The only set left to explore is $\qeq[z]$. This is two parts. If $\qeq[z]$ is empty, then that implies that $c\rint c$ is not a $y$-interval for all rationals $c$. Since singletons $c\rint c$ are a $y$-interval exactly when $c$ is in both $\qleq$ and $\qgeq$, that intersection is empty. Since the elements of $\qeq$ are common to both unions, this implies $\qeq$ is empty. The other part is if $c \in \qeq[z]$. This means $c \yrel c$. But that means that $c$ was an element of the intersection of $\qleq$ and $\qgeq$. Since $\qlt$, $\qeq$, and $\qgt$ are disjoint, this implies $c \in \qeq$ as was needed to be shown. This also shows that, for any given rational $q$, the $q$-cut is mapped to the $q$-relation. 

Thus, given a cut, it maps to a relation that maps to the given cut. The two mappings are inverses to each other.

The above proves 
\begin{theorem}
    The set of rational Dedekind cuts is in canonical bijective correspondence with the set of rational betweenness relations. The representative of the rational $q$ in the relations does correspond to the cut associated with $q$.
\end{theorem}

From this point on, using $x$ in the relation notation $\xrel$ and the cut notation $\cut$, or in the sets of the cut $(\qlt, \qeq, \qgt)$, implies that they are related by this canonical bijection. 

\section{Ordering the Relations}

Dedekind cuts are a complete ordered field as shown in, for example, \cite{rudin}. Rudin uses the lower cuts and this readily translates into this paper's version of Dedekind cuts due to Proposition \ref{br--lowcutdef}. To establish that  the betweenness relations are a complete ordered field, the path will be to show that the natural operations on the betweenness relations is mapped to those operations on their associated cuts. This section will be on the order relations.  

Given two intervals $a \rint b$ and $c\rint d$, the notion that $a \rint b < c\rint d$ is that given any $p$ and $q$ such that $a\rint p\rint b$ and $c\rint q\rint d$, it can be concluded that $p < q$. This includes the endpoints and can be determined by comparing endpoints. Indeed, if $a\leq b$ and $c \leq d$, then $a \rint b < c\rint d$ exactly when $b < c$. Transitivity of interval inequality can readily be shown to hold true. 

From the rational betweenness relation perspective, $x < y$ occurs, by definition, exactly when there exists an $x$-interval $a \rint b$ and a $y$-interval $c\rint d$ such that $a \rint b < c\rint d$. By Proposition \ref{br--different}, given two distinct relations $\xrel$ and $\yrel$, there exists disjoint intervals $a \rint b$ and $c\rint d$ such that $a \xrel b$ and $c \yrel d$. Disjointness means that one of the intervals is wholly greater than the other. By Proposition \ref{br--intersect}, that order relation cannot be reversed. For example, if $a \xrel b$, $c \xrel[y] d$, $u \yrel v$, and $a \rint b < c\rint d$, then $u:v \nless a \rint b$ as $u\rint v \nless c\rint d$ as $u\rint v$ and $c\rint d$ intersect.  To emphasize it as a relation, the inequality may be denoted as $\xrel  \:<\: \yrel$. 

If both relations happen to be  rooted with roots $p$ and $q$, respectively to the relations $\xrel$ and $\yrel$, then $x<y$ if and only if $p < q$ if and only if $p\rint p < q\rint q$. 

For Dedekind cuts, $x < y$ occurs, by definition, exactly when $\qlt \subsetneq \qlt[y]$. This would imply that $\qgt[y] \subset \qgt$. For cuts representing rationals, the solitary elements of $\qeq$ and $\qeq[y]$ will have that ordering though there is no set relation that expresses it. To emphasize it as a cut, the inequality may be denoted as $\cut < \cut[y]$.

The ordering definition of Dedekind cuts is well-established to obey all the desired properties. The task here is to show that the relation ordering and the Dedekind ordering correspond under the bijection which implies the relation ordering also has the same desired properties. 

Assume that $\xrel  \:<\: \yrel$ which implies there exists an $x$-interval $a \rint b$, $a \leq b$,  and a $y$-interval $c\rint d$, $c \leq d$, such that $a \rint b < c\rint d$. Then $a \in \qleq$, $b \in \qgeq$, $c \in \qleq[y]$, $d \in \qgeq[y]$. Since $ b< c$, Lemma \ref{br--cut-unbounded} implies $b \in \qlt[y]$. Since all elements of $\qlt$ are less than $b$, $\qlt \subset \qlt[y]$. Because $b\notin \qlt$ while $b \in \qlt[y]$, it is that case that $\qlt \neq \qlt[y]$. Thus,  $\cut < \cut[y]$. 

For the other direction, assume $\cut < \cut[y]$. By the meaning of the inequality, there exists $b \in \qlt[y]$ which is not in $\qlt$; this implies $b \in \qgeq$. Let $a \in \qlt$ and $d \in \qgt[y]$.  Because there is no greatest element of $\qlt[y]$, let $c\in \qlt[y]$ such that $b < c$. By the definition of the bijection,  $a \xrel b$ and $c \yrel d$. Because $a < b < c < d$, it is the case that $a \rint b < c\rint d$ so that $\xrel  \:<\: \yrel$, as was to be shown. 

\begin{theorem}
For relations $\xrel$ and $\yrel$, $\xrel  \:<\: \yrel$ if and only if $\cut < \cut[y]$.
\end{theorem}


With an ordering, the least upper bound of a set can be queried as to whether it exists or not. Lower cuts give an elegant way to define the least upper bound. Given a non-empty set $E$ of cuts that is bounded above by a cut $z$, then the least upper bound $\alpha$ of $E$ is defined by its lower cut $\qlt[\alpha] = \bigcup_{x \in E} \qlt$. It is straightforward to establish that this is a lower cut, that the cut $\alpha$ is an upper bound, and it is the least upper bound. This is done in detail in the previously cited Rudin \cite{rudin} and establishes that Dedekind cuts are complete. 

 For the relations, the bijection maps a bounded set of relations $\tilde{E}$ to a bounded set of cuts $E$. The least upper bound $\cut[\alpha]$ of $E$ maps to a relation $\xrel[\alpha]$. Because the bijection respects the ordering, the image will be the least upper bound relation of $\tilde{E}$. 

\begin{corollary}
    The set of rational betweenness relations with the interval ordering satisfies the least upper bound property. 
\end{corollary}

By definition of the bijection, $a \xrel[\alpha] b$, $a \leq b$, exactly when $a \in \qleq [\alpha]$ and $b \in \qgeq[\alpha]$. In words, $b$ is a rational upper bound of $E$ while $a$ is a rational lower bound of the set of upper bounds of $E$. This is what one would expect of the least upper bound of the relations. 

A similar bijection dance can be used to establish that the rationals are dense in the relations. But it is rather simple to see it directly as well. For rational betweenness relations, given $\xrel < \yrel$, the task is to show that there is a rational between those two. The inequality implies there exists $a,b,c,d$ such that $a \xrel b$, $c \xrel[y] d$, and $a \rint b \betneq c\rint d$ with $b < c$. Let $m = \frac{b+c}{2}$ be the average of $b$ and $c$. This is a rational number and the relation $\xrel[m]$ satisfies $\xrel  \:<\: \xrel[m]  \:<\: \xrel[y]$ because $a \rint b < m\rint m < c\rint d$ due to $b < m < c$. Thus, the rooted relations are dense in the ordered set of betweenness relations. 

\begin{corollary}
    The set of rooted rational betweenness relations is dense in the set of rational betweenness relations.
\end{corollary}
 

\section{The Arithmetic of Relations}

Having established the canonical bijection between cuts and relations and that the ordered properties work, the final step to do is the arithmetic. As fully proving arithmetic can be a little bit complicated, this paper will take the shortcut of using Dedekind cut arithmetic.

\begin{theorem}
The canonical bijection between cuts and rational betweenness relations is an isomorphism which establishes that the rational betweenness relations form a complete, ordered field with the rational numbers as a dense subset. 
\end{theorem}


The approach will be that the field structure on rational betweenness relations is obtained by transporting the field structure on cuts across the canonical bijection.

\begingroup

\providecommand{\tuple}[2][n]{{}^{#1} #2}


 \noindent\textbf{Convention.} For this proof, if $E_i$ is an indexed expression, then $\tuple{E_i}$ denotes the ordered n-tuple whose $i$-th entry is $E_i$. For example, $\tuple{\qlt[x_i]}$ is a tuple of $n$ lower cuts.

\begin{proof}
    Let $O\cut[<]$ be an arithmetic operator on the space of lower cuts. By \cite{rudin}, pages 17-21, the usual arithmetic operators satisfy the properties of an ordered field. For $n$-ary operators, define $O\cut[](
    \tuple{\cut[x_i]} ) = \overline{O\cut[<]( \tuple{\qlt[x_i]} )}$. In words, take the lower cut of the given cuts, operate on them, then use the completion to cuts via Proposition \ref{br--lowcutdef}. Since the completion to cuts is canonical, equality is preserved under this correspondence.

    Define operators on relations via operators on cuts by using $\xrel[O](\tuple{\!\xrel[x_i]}) = \xrel[z]$ where 
    $\cut[z] = O\cut[](\tuple{\cut[x_i]})$.  As the field properties hold on the arithmetic of the cuts, the properties transfer to the operators on the relations. For example, the commutativity of addition schematically comes from  
    \[
    \begin{aligned}
     \xrel + \xrel[y] &\longrightarrow \cut + \cut[y] \\
     &\longrightarrow \qlt + \qlt[y] \\ 
     &= \qlt[y] + \qlt \\
     &\longrightarrow \cut[y] + \cut \\
     &\longrightarrow \xrel[y] + \xrel   
    \end{aligned}
    \]
     where the equality is established by \cite{rudin}. To do this with the operator notation, let $O(x,y) = x + y$ and $P(x,y)= y + x$. Then the goal is to show $\xrel[O](\xrel, \xrel[y]) = \xrel[P](\xrel, \xrel[y])$ which will follow from $O\cut[<](\qlt[x],\qlt[y]) = P\cut[<] (\qlt[x], \qlt[y])$ as argued below. 
        
    In general, let $O$ and $P$ be two operators such that $O\cut[<](\tuple{\qlt[x_i]}) = P\cut[<](\tuple{\qlt[x_i]})$. Then the goal is to show $\xrel[O](\tuple{\!\xrel[x_i]}) = \xrel[P](\tuple{\!\xrel[x_i]})$ which will hold true, by definition, if and only if $O\cut[](\tuple{\cut[x_i]}) =P\cut[](\tuple{\cut[x_i]}) $, which, in turn, holds if and only if $\overline{O\cut[<](\tuple{\qlt[x_i]})} = \overline{P\cut[<](\tuple{\qlt[x_i]})}$. The last equality follows from  the starting assumption that  $O\cut[<](\tuple{\qlt[x_i]}) = P\cut[<](\tuple{\qlt[x_i]})$.

    In what follows, $\cdot$ will stand in generically for $+$ and $*$. The notation $x^{-1}$ for $\cdot$ will stand in for $-x$ and $\frac{1}{x}$ for $+$ and $*$, respectively. The number $1$ will represent the identity for the generic operator.
    \begin{itemize}
        \item Commutativity. $O(x,y) = x \cdot y$ and $P(x,y)= y \cdot x$
        \item Associativity. $O(x,y,z)= (x\cdot y) \cdot z$ and $P(x,y,z)=x \cdot (y \cdot z)$
        \item Distributive.  $O(x,y,z)= x*(y+x)$ and $P(x,y,z)= (x*y) + (x*z)$
        \item Identity. $O(x) = 1 \cdot x$ and $P(x)= x$
        \item Inverse. $O(x) = x \cdot x^{-1}$ and $P(x)= 1$.
    \end{itemize}
    
    Apply the equality argument above to all of them noting that the lower cut equality for each of them has been established in Rudin. 

    The rest of the theorem was established earlier. 
\end{proof}
\endgroup

While the above suffices, some exploration of what it means for intervals is useful. The arithmetic of the relations  can be described as doing the arithmetic on the endpoints of intervals. This creates a core of intervals which can then be expanded to the full relation. 

To make the correspondence with Dedekind cuts easier, the intervals in the core of the arithmetic will be taken to be Spaced intervals. An $x$-\textbf{Spaced interval} is an $x$-related interval $a \rint b$ such that neither $a$ nor $b$ are roots of the relation. If the labeling is such that $a \leq b$, then it can be concluded that $a < b$, $a \in \qlt$ and $b \in \qgt$.

In what follows, the explicit correspondence of the core intervals to the lower cuts will be explored.  At times, it is necessary to focus on $x > 0$ at which point the notation $\qlt^+$ is useful and it denotes the positive elements of the set of $\qlt$. It is also important to note that having a defined $\qlt^+$ does then yield a unique lower cut, corresponding to $\qlt$, by taking the union with $\qleq[0]$.

\begin{enumerate}
    \item Addition. $u \xrel[x+y] v$ if it is an interval of the form $(a+c)\rint (b+d)$ where $a \rint b$ is an $x$-Spaced interval, $c\rint d$ is a $y$-Spaced interval, $a < b$, and $c < d$. The definition of cut addition is based on $\qlt[x+y] = \{a + c | a \in \qlt, c \in \qlt[y]\} $. This is exactly the lower endpoints of the core addition intervals. 
    \item Negation. $u \xrel[-x] v$ if it is an interval of the form $-b\rint -a$ where $a \rint b$ is an $x$-Spaced interval. The definition of cut negation is based on $\qlt[-x] = \{-b | b \in \qgt\}$. This is exactly the lower endpoints of the core negation intervals. 
    \item Multiplication. $u \xrel[xy] v$  if it is an interval of the form $m\rint M$ where $m = \min(ac, bc, ad, bd)$, $M = \max(ac, bc, ad, bd)$, $a \rint b$ is an $x$-Spaced interval, and $c\rint d$ is a $y$-Spaced interval. For $0 \leq a < b$ and $0 \leq c < d$, this translates to $ac\rint bd$. To correspond to Dedekind cuts, only the positive version is done directly even though the relation version has no need for such a restriction. The other versions are handled by imposing the usual negation and zero multiplication rules. Restricting to the case of $x , y > 0$, the lower cut of $xy$ is based on $\qlt[xy]^+ = \{ac | a \in \qlt^+, c \in \qlt[y]^+ \}$. This is exactly the lower endpoints of the core multiplication intervals of positive numbers. 
    \item Reciprocation. $u \xrel[\frac{1}{x}] v$ if it is an interval of the form $\frac{1}{b}\rint \frac{1}{a}$ where $a \rint b$ is an $x$-Spaced interval which does not contain 0. This is restricted to non-zero numbers. Since one can double negate, with reciprocating in between, it is sufficient to discuss positive cuts. For $x> 0$, the set $\qlt[1/x]^+ = \{\frac{1}{b} | b \in \qgt\}$ which is exactly the set of lower endpoints of the core reciprocation intervals. 
\end{enumerate}


In this approach, the operation of powering to a natural number is taken to be repeated multiplication as that requires no further work. But this can be shown to be equivalent to directly defining the power operation with the core intervals being $a^n \xrel[x^n] b^n$ where $a \rint b$ is an $x$-Spaced interval that does not contain 0. Since the notation is betweenness notation, no ordering of $a^n \rint b^n$ is implied. For even powers, the exclusion of 0 keeps the powering monotone ensuring the endpoints remain endpoints after applying the power. For $x = 0$, $0^n$ is defined as $0$. 


\section{Conclusion}

Rational Betweenness Relations have been established to be a model of real numbers using Dedekind Cuts. The path in this paper is not the only way to establish this. In \cite{taylor24oraclerbr}, the approach is to define another kind of object, called an oracle, which answers Yes or No when given an interval; the intent is that Yes means the real number it represents is in the given interval. It has similar properties to the Rational Betweenness Relations, but allows for a fuzzier, more indeterminate nature. It is closer to numerical work and helps with the actual computational difficulties of dealing with real numbers and the inability to say when two real numbers are definitively equal. 

Using intervals to define real numbers is not new.  A nested sequence of intervals being a real number was proposed by Bachmann not too long after the use of Cauchy sequences and Dedekind cuts. The constructivists, as conveyed by Mark Bridger \cite{bridger}, generalize that approach and use fine families of intervals where are families of intersecting, arbitrarily small intervals to represent real numbers. Both of those approaches lack uniqueness of representatives. The oracles mentioned above are equivalent to the fine families.

There are many definitions of real numbers as surveyed in the excellent overview by Ittay Weiss \cite{ittay-2015}. They have various strengths and weaknesses such as multiple representatives for a given real number or difficulty in handling approximations. Many, such as Dedekind cuts, rely on some particulars of rational numbers that prevent generalization to other spaces. Relations can be generalized to other spaces using containers; the Separation Property is changed to say that given two points in a container, then there exists a subcontainer excluding at least one of the points. Other real number definitions, such as Cauchy sequences, have almost a misleading quality about them, such as having the initial segment of it being completely arbitrary while the infinite tail is not something that one can generally write down. An attempt to tighten that down using tighter bounds on the epsilon condition then leads more to an interval outlook and also tends to complicate arithmetic. 

One particular strength of the relations and related interval definitions, is that they carry definiteness with them naturally. Think about computing a real number for some other computation. The relation approach naturally leads to using intervals. If one has a specific length requirement, the tools and structures required to get to that level come with the relation definition itself. For cuts, it has to be, admittedly rather easily, built up. The cut structure itself seems less relevant. This may happen with definitions, but it is preferable to have definitions that align with the use of the object. 

In particular, imagine trying to get close to $\pi$ and let's say it is known that $\pi$ is between 3 and 4. For the relations, this is asserting that $3 \xrel[\pi] 4$. It is just a restatement of that fact. For the cuts, it would be the assertion that $3 \in \qlt[\pi]$ and $4 \in \qgt[\pi]$. It feels as if the message of the cuts is that we are filling up the sets $\qlt[\pi]$ and $\qgt[\pi]$. It has the feeling of looking away from the real number rather than the relations viewpoint which is to enclose the real number. 

Rational betweenness relations are a new way of perceiving the real numbers. They give a clear and accessible definition while also being able to help with computing out the number. They encourage precise thinking about the imprecision associated with a real number. The concept also has the promise of generalizations to many other spaces. They present a way for how finite beings can meaningfully claim knowledge of an inherently infinite real number. 


\medskip

\normalem %restoring normal emphasis in bibliography 
\printbibliography

\end{document}

